\documentclass[aspectratio=1610, handout]{beamer}%, handout
\usepackage{graphics}
\usepackage{pifont}
\usepackage{ulem}
\usepackage{xcolor}
\usepackage{modifycolor}
\usepackage{lipsum}
\usepackage{fontspec}
\usepackage{caption}
\renewcommand{\figurename}{FIGURE}
\captionsetup[figure]{labelfont={color=leftFootlineColor, scriptsize}, textfont={color=normalBlockColor, scriptsize}}
\usepackage{tabularx}
\newcolumntype{Y}{>{\centering\arraybackslash}X}
\usepackage{tikz}
\usepackage{standalone}
\usepackage{svg}
\usepackage{multicol}
\usepackage{catchfilebetweentags}
\usepackage{xifthen}
\usetikzlibrary{arrows}
\usetikzlibrary{backgrounds}
\setmonofont[
  Contextuals={Alternate}
]{FuraCode Nerd Font}
\setsansfont{Roboto Medium}
\usepackage{pgfpages}


\usepackage[cache=false,outputdir=build]{minted}
\definecolor{bg_code}{HTML}{282828}
\usemintedstyle{darcula}
\setminted{
      fontsize=\scriptsize, 
      linenos,
      numbersep=0pt,
      gobble=5,
      framesep=3mm} 
      \renewcommand{\theFancyVerbLine}{\texttt{{{\arabic{FancyVerbLine}}}}}
\usepackage{adjustbox}
\usepackage{environ}%
\usepackage{tikz}
\usetikzlibrary{patterns}
\usepackage[french]{babel}
\usepackage{beamerthemesideblue}


\setbeamerfont{note page}{size=\footnotesize}
\setbeamertemplate{note page}[custom]
\setbeamercolor{note page}{bg=backgroundColor, fg=white}

\newcommand{\insertlicense}{\includegraphics[height=1cm]{img/by-sa.png}}
\title[]{La rétroingénierie appliquée à Android}
\subtitle{La traque aux traqueurs}
\author{Maxime Catrice}
\date{\today}
\titlegraphic{\includegraphics[height=0.95cm]{logo.png}\bigbreak\insertlicense}

\setbeamertemplate{title page}[default][colsep=-4bp,rounded=true,shadow=false]
\setbeamercolor{background canvas}{bg=backgroundColor}
\setbeamertemplate{blocks}[rounded][shadow=false]
\beamertemplatenavigationsymbolsempty
\newcounter{acolumn}%  Number of current column
\newlength{\acolumnmaxheight}%   Maximum column height
%%%%%%%%%%%%%%%%%%%%%%%%%


\makeatletter

% `column` replacement to measure height
\newenvironment{@acolumn}[1]{%
    \stepcounter{acolumn}%
    \begin{lrbox}{\@tempboxa}%
    \begin{minipage}{#1}%
}{%
    \end{minipage}
    \end{lrbox}
    \@tempdimc=\dimexpr\ht\@tempboxa+\dp\@tempboxa\relax
    % Save height of this column:
    \expandafter\xdef\csname acolumn@height@\roman{acolumn}\endcsname{\the\@tempdimc}%
    % Save maximum height
    \ifdim\@tempdimc>\acolumnmaxheight
        \global\acolumnmaxheight=\@tempdimc
    \fi
}

% `column` wrapper which sets the height beforehand
\newenvironment{@@acolumn}[1]{%
    \stepcounter{acolumn}%
    % The \autoheight macro contains a \vspace macro with the maximum height minus the natural column height
    \edef\autoheight{\noexpand\vspace*{\dimexpr\acolumnmaxheight-\csname acolumn@height@\roman{acolumn}\endcsname\relax}}%
    % Call original `column`:
    \orig@column{#1}%
}{%
    \endorig@column
}

% Save orignal `column` environment away
\let\orig@column\column
\let\endorig@column\endcolumn

% `columns` variant with automatic height adjustment
\NewEnviron{acolumns}[1][]{%
    % Init vars:
    \setcounter{acolumn}{0}%
    \setlength{\acolumnmaxheight}{0pt}%
    \def\autoheight{\vspace*{0pt}}%
    % Set `column` environment to special measuring environment
    \let\column\@acolumn
    \let\endcolumn\end@acolumn
    \BODY% measure heights
    % Reset counter for second processing round
    \setcounter{acolumn}{0}%
    % Set `column` environment to wrapper
    \let\column\@@acolumn
    \let\endcolumn\end@@acolumn
    % Finally process columns now for real
    \begin{columns}[#1]%
        \BODY
    \end{columns}%
}
\makeatother
%%%%%%%%%%%%%%%%%%%%%%%%%

\NewEnviron{frameNoSB}{
\makeatletter
\setbeamertemplate{sidebar canvas left}{}
\setbeamertemplate{sidebar left}{}
\makeatother
\begin{frame}
    \BODY
%\tableofcontents
\end{frame}
}

\NewEnviron{frameTitle}{
\setbeamertemplate{frametitle}[default][center]
\makeatletter
\setbeamertemplate{headline}{\color{backgroundColor}45\newline 45\newline 45\newline 45\newline 45\newline 45\newline 45\newline 45\newline 45\newline 45}
\setbeamertemplate{sidebar canvas left}{}
\setbeamertemplate{sidebar left}{}



\makeatother
\begin{frame}
\begin{minipage}[c]{\linewidth-\beamerleftmargin+\beamerrightmargin}
\BODY
\end{minipage}
\end{frame}
}

\setbeamerfont{title}{series=\bfseries,parent=structure}
\setbeamerfont{subtitle}{size=\Large,series=,parent=structure}

\makeatletter
\newlength\beamerleftmargin
\setlength\beamerleftmargin{\Gm@lmargin}

\newlength\beamerrightmargin
\setlength\beamerrightmargin{\Gm@rmargin}
\makeatother

\newcommand{\nologo}{\setbeamertemplate{logo}{}}

\newcommand{\slidetitle}[1][]{
  \frametitle{\insertsection\ifthenelse{\equal{#1}{}}{}{: #1}}
}
\newcommand{\mono}[1]{
\texttt{#1}
}

\setbeameroption{show notes on second screen=top}
\begin{document}
 
\section{L'analyse réseau}
\begin{frame}[t]
  \slidetitle[]
  \note{\onslide<1->{\ExecuteMetaData[007-anares_notes.tex]{Anares}}}
  \note{\onslide<2->{\ExecuteMetaData[007-anares_notes.tex]{Objectifs}}}
  \note{\onslide<7->{\ExecuteMetaData[007-anares_notes.tex]{Environnement}}}
  \begin{block}{Qu'est ce que l'analyse réseau?}<2->
    \onslide<3->{
      Intercepter le traffic entrant et sortant de l'application, pour détermier qui sont les destinataires et comment sont échangés les messages
    }
    \end{block}
  \begin{columns}
    \begin{column}{0.7\linewidth}
      \begin{block}{Objectifs:}<4->
        \begin{itemize}
        \item<5-> Déterminer les échanges effectués par l'application
        \item<6-> Lire le traffic http
        \item<7-> Déchiffrer le traffic https
        \end{itemize}
      \end{block}
      \begin{block}{Environnement utilisé}<7->
        \begin{itemize}
          \item<8-> Emulateur genymotion avec ProxyDroid
          \item<9-> WireShark (Analyseur de paquet)
          \item<10-> Xposed: JustTrustMe
          \end{itemize}
      \end{block}    
    \end{column}
    \begin{column}{0.25\linewidth}
      \centering
      \includegraphics[width=0.75\linewidth]{img/android_network.pdf}
    \end{column}
  \end{columns}
  \vfill
\end{frame}

\begin{frame}[t]
  \slidetitle[Principe]
  \note{\onslide<7->{\ExecuteMetaData[007-anares_notes.tex]{MITM}}}
  \begin{figure}
    \vspace{-0.65cm}
    \documentclass[aspectratio=1610, handout]{beamer}%, handout
    \usepackage{graphics}
\usepackage{pifont}
\usepackage{ulem}
\usepackage{xcolor}
\usepackage{modifycolor}
\usepackage{lipsum}
\usepackage{fontspec}
\usepackage{caption}
\renewcommand{\figurename}{FIGURE}
\captionsetup[figure]{labelfont={color=leftFootlineColor, scriptsize}, textfont={color=normalBlockColor, scriptsize}}
\usepackage{tabularx}
\newcolumntype{Y}{>{\centering\arraybackslash}X}
\usepackage{tikz}
\usepackage{standalone}
\usepackage{svg}
\usepackage{multicol}
\usepackage{catchfilebetweentags}
\usepackage{xifthen}
\usetikzlibrary{arrows}
\usetikzlibrary{backgrounds}
\setmonofont[
  Contextuals={Alternate}
]{FuraCode Nerd Font}
\setsansfont{Roboto Medium}
\usepackage{pgfpages}


\usepackage[cache=false,outputdir=build]{minted}
\definecolor{bg_code}{HTML}{282828}
\usemintedstyle{darcula}
\setminted{
      fontsize=\scriptsize, 
      linenos,
      numbersep=0pt,
      gobble=5,
      framesep=3mm} 
      \renewcommand{\theFancyVerbLine}{\texttt{{{\arabic{FancyVerbLine}}}}}
\usepackage{adjustbox}
\usepackage{environ}%
\usepackage{tikz}
\usetikzlibrary{patterns}
\usepackage[french]{babel}
\usepackage{beamerthemesideblue}


\setbeamerfont{note page}{size=\footnotesize}
\setbeamertemplate{note page}[custom]
\setbeamercolor{note page}{bg=backgroundColor, fg=white}

\newcommand{\insertlicense}{\includegraphics[height=1cm]{img/by-sa.png}}
\title[]{La rétroingénierie appliquée à Android}
\subtitle{La traque aux traqueurs}
\author{Maxime Catrice}
\date{\today}
\titlegraphic{\includegraphics[height=0.95cm]{logo.png}\bigbreak\insertlicense}

\setbeamertemplate{title page}[default][colsep=-4bp,rounded=true,shadow=false]
\setbeamercolor{background canvas}{bg=backgroundColor}
\setbeamertemplate{blocks}[rounded][shadow=false]
\beamertemplatenavigationsymbolsempty
\newcounter{acolumn}%  Number of current column
\newlength{\acolumnmaxheight}%   Maximum column height
%%%%%%%%%%%%%%%%%%%%%%%%%


\makeatletter

% `column` replacement to measure height
\newenvironment{@acolumn}[1]{%
    \stepcounter{acolumn}%
    \begin{lrbox}{\@tempboxa}%
    \begin{minipage}{#1}%
}{%
    \end{minipage}
    \end{lrbox}
    \@tempdimc=\dimexpr\ht\@tempboxa+\dp\@tempboxa\relax
    % Save height of this column:
    \expandafter\xdef\csname acolumn@height@\roman{acolumn}\endcsname{\the\@tempdimc}%
    % Save maximum height
    \ifdim\@tempdimc>\acolumnmaxheight
        \global\acolumnmaxheight=\@tempdimc
    \fi
}

% `column` wrapper which sets the height beforehand
\newenvironment{@@acolumn}[1]{%
    \stepcounter{acolumn}%
    % The \autoheight macro contains a \vspace macro with the maximum height minus the natural column height
    \edef\autoheight{\noexpand\vspace*{\dimexpr\acolumnmaxheight-\csname acolumn@height@\roman{acolumn}\endcsname\relax}}%
    % Call original `column`:
    \orig@column{#1}%
}{%
    \endorig@column
}

% Save orignal `column` environment away
\let\orig@column\column
\let\endorig@column\endcolumn

% `columns` variant with automatic height adjustment
\NewEnviron{acolumns}[1][]{%
    % Init vars:
    \setcounter{acolumn}{0}%
    \setlength{\acolumnmaxheight}{0pt}%
    \def\autoheight{\vspace*{0pt}}%
    % Set `column` environment to special measuring environment
    \let\column\@acolumn
    \let\endcolumn\end@acolumn
    \BODY% measure heights
    % Reset counter for second processing round
    \setcounter{acolumn}{0}%
    % Set `column` environment to wrapper
    \let\column\@@acolumn
    \let\endcolumn\end@@acolumn
    % Finally process columns now for real
    \begin{columns}[#1]%
        \BODY
    \end{columns}%
}
\makeatother
%%%%%%%%%%%%%%%%%%%%%%%%%

\NewEnviron{frameNoSB}{
\makeatletter
\setbeamertemplate{sidebar canvas left}{}
\setbeamertemplate{sidebar left}{}
\makeatother
\begin{frame}
    \BODY
%\tableofcontents
\end{frame}
}

\NewEnviron{frameTitle}{
\setbeamertemplate{frametitle}[default][center]
\makeatletter
\setbeamertemplate{headline}{\color{backgroundColor}45\newline 45\newline 45\newline 45\newline 45\newline 45\newline 45\newline 45\newline 45\newline 45}
\setbeamertemplate{sidebar canvas left}{}
\setbeamertemplate{sidebar left}{}



\makeatother
\begin{frame}
\begin{minipage}[c]{\linewidth-\beamerleftmargin+\beamerrightmargin}
\BODY
\end{minipage}
\end{frame}
}

\setbeamerfont{title}{series=\bfseries,parent=structure}
\setbeamerfont{subtitle}{size=\Large,series=,parent=structure}

\makeatletter
\newlength\beamerleftmargin
\setlength\beamerleftmargin{\Gm@lmargin}

\newlength\beamerrightmargin
\setlength\beamerrightmargin{\Gm@rmargin}
\makeatother

\newcommand{\nologo}{\setbeamertemplate{logo}{}}

\newcommand{\slidetitle}[1][]{
  \frametitle{\insertsection\ifthenelse{\equal{#1}{}}{}{: #1}}
}
\newcommand{\mono}[1]{
\texttt{#1}
}

    \definecolor{backgroundColor}{HTML}{2D2D2D}
    \usetikzlibrary{arrows}
    \usetikzlibrary{backgrounds}
    \begin{document}
    \begin{tikzpicture}[scale=0.45, every node/.style={scale=0.6},color=white,]
        \tikzstyle{myarrows}=[thick, line width=1.7mm,-open triangle 90,
        postaction={draw, line width=6.5mm, shorten >=4mm, -},
        postaction={draw, backgroundColor, shorten >= 1.3mm, -triangle 90, line width=1mm},
        postaction={draw, line width=5mm, backgroundColor, shorten >= 4mm, -}]
        
    \node at (0,0) {
    \includegraphics[height=3cm]{img/smartphone.pdf}
    };
    \node at (0,-2.5) {
        \LARGE{Client}
    };

    \onslide<1-1|handout:1>{
        \node[xscale=-1] at (2.5,2) {
            \includegraphics[width=1.5cm]{img/key_green.pdf}
        };
        \node at (20,-2) {
            \includegraphics[width=1.5cm]{img/key_green.pdf}
        };
    }
    \node at (24,0) {
        \includegraphics[height=3cm]{img/server.pdf}
    };
    \node at (24,-2.5) {
        \LARGE{Serveur}
    };
    \onslide<2-3|handout:2>{
    \node[xscale=-1] at (2.5,2) {
        \includegraphics[width=1.5cm]{img/key_purple.pdf}
    };
    \node[xscale=1] at (8.5,-2) {
        \includegraphics[width=1.5cm]{img/key_purple.pdf}
    };
    \node[xscale=-1] at (15.5,2) {
        \includegraphics[width=1.5cm]{img/key_orange.pdf}
    };
    \node at (20,-2) {
        \includegraphics[width=1.5cm]{img/key_orange.pdf}
    };
    %\onslide<2-2>{
    \node at (12,0) {
        \includegraphics[height=3cm]{img/computer.pdf}
    };
    \node at (12,-2.5) {
        \LARGE{MITM}
    };
    }
    \node[inner sep=0,minimum size=0] at (1.5,1) (C1) {}; % invisible node
    \node[inner sep=0,minimum size=0] at (21,1) (S1) {}; % invisible node

    \node[inner sep=0,minimum size=0] at (1.5,-1) (C2) {}; % invisible node
    \node[inner sep=0,minimum size=0] at (21,-1) (S2) {}; % invisible node

    
    \node[inner sep=0,minimum size=0] at (9.5,1) (ML1) {}; % invisible node
    \node[inner sep=0,minimum size=0] at (9.5,-1) (ML2) {}; % invisible node
    \node[inner sep=0,minimum size=0] at (14.5,1) (MR1) {}; % invisible node
    \node[inner sep=0,minimum size=0] at (14.5,-1) (MR2) {}; % invisible node

    \node[inner sep=0,minimum size=0] at (12,-3) (MB) {}; % invisible node
    \node[inner sep=0,minimum size=0] at (12,-10) (WT) {}; % invisible node

    \onslide<1-1|handout:1>{
        \draw [-triangle 60, line width=3pt, green] (C1) -- (S1);
        \draw [-triangle 60, line width=3pt, green] (S2) -- (C2);
    }
    \onslide<2-3|handout:2>{
        \draw [-triangle 60, line width=3pt, orange] (S2) -- (MR2);
        \draw [-triangle 60, line width=3pt, orange] (MR1) -- (S1);

        \draw [-triangle 60, line width=3pt, purple] (C1) -- (ML1);
        \draw [-triangle 60, line width=3pt, purple] (ML2) -- (C2);
    }

    \onslide<3-3|handout:2>{
        \draw [-triangle 60, line width=3pt, normalBlockColor] (MB) -- (WT);

        \node at (12,-12) {
            \includegraphics[height=3cm]{img/wireshark.pdf}
        };
    }
    \end{tikzpicture}
\end{document}
    \caption{Principe d'une attaque man-in-the-middle}
  \end{figure}
\end{frame}
\end{document}
